% !TEX root = ../mainfile.tex
% !TEX encoding = UTF-8 Unicode
% !TEX spellcheck = en_US

\chapter{Components of the Work}

The structure of a scientific work typically follows a fixed format and always includes a title page, a table of contents, the main text, a bibliography, and a declaration of academic honesty. Optionally, abbreviation, symbol, table, and figure lists, as well as an appendix (or multiple appendices), can be added. The individual elements are briefly explained below.

\section{Title Page}
The title page for bachelor's theses has been designed according to the guidelines of the Department of Economics. Please insert the missing information into the document \texttt{title.tex}. Deviations are possible, but the listed information should not be omitted even with a different design. For seminar papers, the title page should be modified accordingly. In particular, the title of the seminar and the lecturer should be specified.

\section{Table of Contents}
The table of contents updates automatically. To do this, \LaTeX\ requires setting twice – you must therefore call the \LaTeX\ program at least twice.

\section{List of Figures and Tables}
Unlike the table of contents, lists of figures and tables are not necessary for every work. They are useful when working with many figures and tables. In this case, figures and tables must be numbered separately (this happens automatically). Both lists are automatically updated by \LaTeX\, provided you include appropriate captions for each figure and table (see Section \ref{sec:floats}).

\section{List of Abbreviations and Symbols}
Like the lists of figures and tables, lists of abbreviations and symbols are not mandatory. In the list of abbreviations, include abbreviations that are not commonly used in everyday language. Abbreviations like DDR, PC, or e.g. do not need to be listed. In the symbol list, include all symbols you use. Make sure to explain abbreviations and symbols the first time they are used.

\section{Main Text with Introduction, Main Part, and Conclusion}

This is the main part of the work, where you answer your chosen research question. This part is divided into at least three chapters: "Introduction," "Main Part," and "Conclusion," with the main part usually divided into several chapters and subchapters. For further information on the content design of your scientific work, please refer to the "Tips for Scientific Papers," which you can download or read online on the website of the \href{https://www.vwl.uni-mannheim.de/studium/bachelorstudium/schreibberatung/weiterfuehrendes-material-und-vorlagen/}{\color{blue} Writing Advisory Service (Schreibberatung)}.

\section{Bibliography}
In the bibliography, list all sources you have used. The sources should be sorted alphabetically by the author's last name. Please adhere to Harvard citation style and be consistent. You can also find detailed information on this in the "Tips for Scientific Papers."

\section{Appendix}
The appendix contains additional information referenced in the main text but would take up too much space there. Typically, this includes questionnaires, tables, statistics, data analyses, or transcripts of interviews. However, please note that all essential information for answering your research question should be included in the main part. Therefore, you must not pack extensive important text passages or figures into the appendix to circumvent page restrictions.

\section{Declaration of Academic Honesty}
With this declaration, you affirm that you have independently authored your scientific work and have cited all sources used. Deception can lead to revocation of your bachelor's degree. The declaration must be kept in its original wording, manually signed, and dated on the copies to be submitted.